\section{Dashboard de observabilidad}

El dashboard es la cara visible del cluster: una sola página web que monitorea
los tres nodos a la vez y, para cada uno, pinta los datos de negocio (cuentas,
saldo, transferencias, secuencia de la última transacción) junto con las
métricas de infraestructura del sistema operativo de su máquina (CPU, RAM y
Disco). El armado descansa en tres endpoints HTTP, registrados en
\codigo{app/Node.java} sobre el router \codigo{net/Routes.java}: \codigo{GET /}
sirve la página, \codigo{GET /panel} agrega la vista de todo el cluster y
\codigo{GET /api/stats} expone las métricas de un solo nodo. El extra de la
rúbrica (auto-refresh, sin tener que pulsar ``Refrescar'') se cubre con un
\codigo{setInterval} de polling en el JavaScript de la página.

\subsection{La interfaz web}

La página vive como recurso estático en \codigo{src/main/resources/dashboard.html}
y se entrega tal cual desde el classpath por \codigo{api/DashboardEndpoint.java}:
ese endpoint solo valida que el método sea \codigo{GET} (responde \codigo{405} si
no) y vuelca el HTML con \codigo{Reply.html(200, ...)}, leyendo el recurso
\codigo{/dashboard.html} en cada petición. No tiene colaboradores ni lógica de
negocio; toda la inteligencia está en el JavaScript del propio documento.

Es una SPA en HTML/CSS/JS puro, sin dependencias ni librerías externas. Al
cargar, y cada vez que se refresca, la función \codigo{loadPanel()} hace
\codigo{fetch("/panel")} pidiendo JSON y, con la respuesta, repinta tres zonas:
el \emph{hero} (saldo total del cluster, número de cuentas y última transacción,
tomados del nodo líder), un sello de consistencia que compara los saldos totales
de los nodos en línea para indicar si el invariante se conserva, y una tarjeta
por nodo. Cada tarjeta dibuja el rol (\codigo{Lider}/\codigo{Replica}), el estado
(\codigo{Activo}/\codigo{Caido}), el saldo del nodo, los cuatro contadores de
negocio (cuentas, transferencias, última tx y transacciones en Cloud Storage) y
tres barras (\emph{gauges}) para CPU, RAM y Disco. Si \codigo{/panel} falla, se
muestra un \emph{banner} de error en vez de romper la página.

El auto-refresh, que da los 15 puntos extra, es un \emph{toggle} ``En vivo'' que
\textbf{arranca encendido por defecto}: el checkbox lleva \codigo{checked} y, al
cargar, la página llama \codigo{toggleReactive(checkbox.checked)}, de modo que el
tablero empieza a refrescarse solo \emph{sin pulsar ``Refrescar''}, que es justo lo
que pide el extra. \codigo{toggleReactive()} arranca un \codigo{setInterval} que
vuelve a llamar a \codigo{loadPanel()} cada \codigo{REACTIVE\_INTERVAL\_MS}
(1500 ms); al desactivar el toggle se limpia el temporizador con
\codigo{clearInterval}.

\begin{lstlisting}[language=none,caption={Polling automatico del dashboard (dashboard.html)}]
// Poll cadence (ms) used when the reactive toggle is enabled.
const REACTIVE_INTERVAL_MS = 1500;
let timer = null;
// ...
function toggleReactive(enabled) {
  if (timer) { clearInterval(timer); timer = null; }
  if (enabled) timer = setInterval(loadPanel, REACTIVE_INTERVAL_MS);
  loadPanel();
}
\end{lstlisting}

\subsection{Datos de negocio}

Los datos de negocio de cada nodo nacen en \codigo{app/NodeStats.java}, un
\emph{snapshot} en vivo que reparte las cifras entre dos colaboradores: del
\codigo{core/Ledger.java} salen el número de cuentas (\codigo{accountCount()} vía
\codigo{ledger.size()}) y el saldo total en centavos (\codigo{totalCents()}); del
\codigo{core/Bank.java} salen el número de transferencias confirmadas
(\codigo{transferCount()} vía \codigo{bank.appliedCount()}) y la secuencia de la
última transacción aplicada (\codigo{lastTxId()} vía \codigo{bank.lastSeq()}). El
método \codigo{toJson()} serializa todo a un \codigo{JsonObject} cuyas claves
coinciden exactamente con los campos que consume el JavaScript: \codigo{nodeId},
\codigo{role}, \codigo{status}, \codigo{accountCount}, \codigo{totalBalance}
(convertido a decimal con \codigo{Money.toDecimal()}), \codigo{transferCount},
\codigo{lastTxId}, los tres porcentajes y \codigo{gcsCount}.

Quien expone ese JSON de un solo nodo es \codigo{api/StatsEndpoint.java} en la
ruta \codigo{GET /api/stats}: valida el método y devuelve
\codigo{Reply.json(200, stats.toJson())}. Sobre él se monta
\codigo{api/PanelEndpoint.java} (ruta \codigo{GET /panel}), que arma la vista del
cluster completo: parte del \codigo{toJson()} del nodo local y, en paralelo,
consulta el \codigo{/api/stats} de cada peer listado en \codigo{TES\_PEERS}
(\codigo{NodeConfig}). Las peticiones se lanzan concurrentes con
\codigo{CompletableFuture} para que un nodo caído sume a lo sumo un \emph{timeout}
a la latencia del panel, no uno por peer; un peer inalcanzable no rompe la
respuesta, sino que se reporta como nodo en estado \codigo{down}. El resultado es
un único documento \codigo{\{"nodes":[...]\}} con una entrada por nodo.

\begin{lstlisting}[language=Java,caption={PanelEndpoint agrega el nodo local y los peers}]
StringBuilder nodes = new StringBuilder();
nodes.append("{\"nodes\":[").append(stats.toJson());
for (CompletableFuture<String> peer : pending) {
    nodes.append(',').append(peer.join());
}
nodes.append("]}");
return Reply.json(200, nodes.toString());
\end{lstlisting}

\subsection{Métricas de infraestructura del SO}

Las tres barras de CPU, RAM y Disco no son inventadas: \codigo{NodeStats} las lee
del sistema operativo de cada VM con Java puro, sin librerías externas. El CPU se
toma de \codigo{/proc/stat}: \codigo{cpuPercent()} suma los \emph{jiffies} de la
línea agregada \codigo{cpu} y calcula la fracción ocupada como
$1 - \Delta\text{idle}/\Delta\text{total}$ entre dos lecturas consecutivas (entre
dos refrescos del dashboard), guardando el muestreo previo bajo
\codigo{cpuLock} para que las peticiones concurrentes no se pisen. La RAM viene
de \codigo{/proc/meminfo} (\codigo{ramPercent()} compara \codigo{MemTotal} contra
\codigo{MemAvailable}). El Disco se calcula con \codigo{java.io.File} sobre la
raíz \codigo{/}: \codigo{diskPercent()} usa \codigo{getTotalSpace()} y
\codigo{getUsableSpace()}. Las tres pasan por \codigo{clamp()} al rango
$[0, 100]$ y por \codigo{round1()} a un decimal antes de serializarse.

\begin{lstlisting}[language=Java,caption={Lectura de \%RAM y \%Disco del SO (NodeStats)}]
public double ramPercent() {
    try {
        long total = readMeminfo("MemTotal");
        long available = readMeminfo("MemAvailable");
        if (total <= 0L) {
            return 0.0;
        }
        return clamp((double) (total - available) / total * 100.0);
    } catch (RuntimeException | IOException e) {
        return 0.0;
    }
}
\end{lstlisting}

Estos porcentajes viajan al tablero por la misma vía que los datos de negocio:
\codigo{toJson()} los empaqueta como \codigo{cpuPercent}, \codigo{ramPercent} y
\codigo{diskPercent}; \codigo{StatsEndpoint} los publica en \codigo{/api/stats},
\codigo{PanelEndpoint} los agrega en \codigo{/panel}, y el JavaScript de
\codigo{dashboard.html} los dibuja como \emph{gauges} (la barra se pinta en tono
de alarma cuando el valor llega al 80\%). Así, cada métrica del cuadro proviene
directamente del SO de la máquina que la reporta.
